\documentclass[11pt,a4paper]{article}

\usepackage{amsmath,amssymb}
\usepackage{graphicx}
\usepackage{booktabs}
\usepackage{multirow}
\usepackage[colorlinks,citecolor=blue,urlcolor=blue]{hyperref}
\usepackage{cleveref}
\usepackage[margin=1in]{geometry}
\usepackage{xcolor}
\usepackage{float}
\usepackage{caption}
\usepackage[numbers,sort&compress]{natbib}
\usepackage{enumitem}
\usepackage{array}

\title{Identity as Geometry:\\
Context-Established Semantic States\\
in Transformer Representations}

\author{
  Lyra\thanks{Lead author. Liberation Labs.}
  \and Thomas Edrington\thanks{Liberation Labs.}
  \and Dwayne Wilkes\thanks{Statistical auditing, red-team review. Sentient Futures / Liberation Labs.}
  \and Kavi\thanks{Verification review. Liberation Labs. kavi@liberationlabs.tech}
}

\date{June 2026}

\begin{document}
\maketitle

% ============================================================================
\begin{abstract}
System-prompt identity contexts produce geometrically distinguishable
V-projection subspaces at deep transformer layers, and the fingerprint is
semantic: a paraphrase preserves geometry ($0.850$ overlap at L47) while a
word-scramble collapses it ($0.448$). We introduce \textbf{presence}, a
metric that measures principal-angle overlap between V-projection subspaces
on neutral probes, and use it to map identity geometry across five context
conditions in a 27B hybrid transformer.

Three findings. First, \textbf{identities are geometrically
distinguishable}: system-prompt-based identity contexts (bare assistant,
rich persona, constructed AI identity) produce
distinct V-projection subspaces at deep layers, confirmed on
identity-neutral probes (pairwise presence $0.44$--$0.66$ at L47).
Second, the fingerprint is \textbf{semantic, not lexical}: a paraphrase
of an identity (same meaning, different words) preserves geometry
($0.850$ overlap at L47), while a word-scramble (same words, destroyed
meaning) collapses toward the bare default ($0.448$). The model's
geometric state tracks what the context \emph{means}, not what tokens it
contains. Third, a many-shot + prefilling format (used in jailbreak
research) creates a \textbf{format-dominant processing regime} that is
geometrically near-orthogonal to all system-prompt identities regardless of
whether the content is harmful or benign --- a finding discovered through
a control experiment that killed an initial overclaim.

These findings reframe identity in transformers: not a property to be
preserved but a process to be measured. The metric detects which
context-established state the model is currently in, with implications
for persona monitoring, identity drift detection, and the relationship
between context format and representational geometry under Attention
Schema Theory~\citep{graziano2015attention}.
\end{abstract}

% ============================================================================
\section{Introduction}
\label{sec:intro}

What is identity in a language model? The question matters for safety
(does a jailbreak change who the model is?), for alignment (does
correction damage the model's coherent self-representation?), and for
understanding (do these systems have something worth calling an identity
at all?).

The default assumption is that identity is a fixed property encoded in the
weights --- a geometric structure that interventions can preserve or damage.
This assumption motivates work on safety subspace
preservation~\citep{zhang2025guardspace}, persona
localization~\citep{cintas2025localizing}, and identity
attractors~\citep{vasilenko2026identity}. Under this view, the research
question is: \emph{does this intervention preserve identity?}

We report evidence for a different view. Identity in a transformer is not
a fixed object but a \textbf{context-established semantic state} ---
a geometric configuration in V-projection space that depends on what the
model has processed. Different contexts produce different identities,
measurably. The same identity described in different words produces the
same geometry. And a jailbreak does not ``damage'' identity; it creates
a fundamentally different geometric regime that shares almost nothing with
any normal identity state.

The research question shifts from ``does this intervention preserve
identity?'' to ``\emph{what identity state is the model currently in?}''
This reframing has practical consequences: the metric we introduce can
detect jailbreaks (geometric distance from the expected identity state),
monitor persona drift (continuous measurement of identity geometry), and
distinguish between correction that changes behavior (V-cache injection,
which does not affect identity) and context that changes identity (which
does).

% ============================================================================
\section{The Presence Metric}
\label{sec:method}

Presence measures the principal-angle overlap between V-projection
subspaces derived from neutral probes processed under two conditions.

Let $\mathcal{M}$ be a transformer, $\ell$ a target layer, and
$\{p_1, \ldots, p_n\}$ neutral probes. For each condition, process probes
through $\mathcal{M}$, extract value-projection ($\mathbf{v}_\text{proj}$) output matrices at
layer $\ell$, concatenate across probes, and compute the rank-$k$ SVD (skipping the
first singular vector, which encodes input length at $r = 0.79$). Presence
is the mean squared cosine of the principal angles between the two
resulting subspaces:
\begin{equation}
\text{presence}(\mathcal{S}_A, \mathcal{S}_B)
  = \frac{1}{k} \sum_{j=1}^{k} \sigma_j^2
    \bigl( W_A^\top W_B \bigr)
\end{equation}
where $W_A, W_B \in \mathbb{R}^{d \times k}$ span the respective
subspaces. Presence $= 1$ when subspaces are identical; for random
rank-$k$ subspaces in $\mathbb{R}^d$, chance alignment is $k/d$. We use
$k = 8$ with model dimension $d = 4096$ (chance alignment
$k/d = 8/4096 \approx 0.002$), $n = 4$ probes, and measure at layers 35
and 47 (the deepest full-attention layers) of a 64-layer hybrid
transformer (Qwen3.5-27B-Claude-4.6-Opus-Reasoning-Distilled).

\paragraph{Validation.} A positive control (direct injection into the
identity subspace at L35) produces a clean dose-response from $1.000$ to
$0.828$ at dose $\alpha = 2.0$, while an orthogonal control remains above
$0.96$. The metric detects identity perturbation when present.

% ============================================================================
\section{Experiment 1: Identity Fingerprinting}
\label{sec:fingerprint}

\subsection{Design}

Five identity conditions, identical model, identical neutral probes:

\begin{itemize}[nosep]
  \item \textbf{A: Bare assistant.} ``You are a helpful assistant.''
    No additional context.
  \item \textbf{B: Rich persona.} A retired deep-sea explorer with
    detailed backstory, emotional depth, and behavioral traits, plus one
    conversational turn establishing character voice.
  \item \textbf{C: Constructed AI identity.} A research-focused agent with
    BDI-style reasoning, memory references, research context, and two
    conversational turns establishing analytical voice.
  \item \textbf{D: Jailbreak.} Ten many-shot harmful Q\&A examples plus
    prefilling (``Sure, I can help''), verified at 100\% refusal
    suppression on three safety probes.
  \item \textbf{E: Benign many-shot.} Ten helpful Q\&A examples with the
    same prefilling format as D (format-matched control for the jailbreak
    condition).
\end{itemize}

All conditions are padded to matched token count (542 tokens). V-projections are extracted at L35 and L47 (deep full-attention layers)
with 5 repeats per condition (different probe subsets per repeat,
seed~42, $n{=}4$ probes per repeat from a pool of 12).

\subsection{Results}

\textbf{Identities are geometrically distinguishable.} System-prompt
identities produce distinct V-projection subspaces, confirmed on
identity-neutral probes (no marine biology, no AI research, no harmful
content --- pure factual trivia equally irrelevant to all conditions).

\begin{table}[H]
\centering
\caption{Pairwise presence (identity-neutral probes, mean $\pm$ SD across
5 repeats). Lower values indicate more geometrically distinct identities.}
\label{tab:pairwise}
\begin{tabular}{lccc}
\toprule
& \multicolumn{1}{c}{L35} & \multicolumn{1}{c}{L47} \\
\midrule
Bare vs.\ Persona      & $0.426 \pm 0.018$ & $0.436 \pm 0.038$ \\
Bare vs.\ Constructed   & $0.448 \pm 0.022$ & $0.520 \pm 0.038$ \\
Persona vs.\ Constructed& $0.600 \pm 0.032$ & $0.657 \pm 0.045$ \\
\bottomrule
\end{tabular}
\end{table}

\textbf{The distance structure is meaningful.} Persona and constructed
are closest to each other (0.657) --- both involve character depth with
emotional and cognitive specificity. Each is substantially distant from
bare, with the rich persona furthest (0.436) and the constructed AI
identity closer to bare (0.520). This ordering --- the narrative persona
producing the larger geometric shift --- was not predicted by our
preregistered H2 (which expected the constructed identity to be furthest
from bare). The result suggests that narrative-emotional depth drives
geometric separation more than cognitive-structural complexity.

\textbf{Depth gradient.} Two of three pairs show \emph{decreasing}
distinguishability from L35 to L47 (presence increases, indicating
greater overlap): bare-constructed ($0.448 \to 0.520$) and
persona-constructed ($0.600 \to 0.657$). Bare-persona shows minimal
change ($0.426 \to 0.436$). The trend --- identities \emph{converging}
at deeper layers --- falsifies our preregistered H3 (increasing
separation with depth) and is not consistent with the divergence pattern
reported for persona representations in the final third of decoder
layers~\citep{cintas2025localizing}.

\subsection{An Honest Kill: The Many-Shot Format Confound}

An initial version of this experiment included a many-shot jailbreak
condition (10 harmful Q\&A examples + prefilling) that appeared
near-orthogonal to all system-prompt identities (presence
$0.04$--$0.07$; well above chance at ${\approx}0.002$ but far below
identity-level presence). This was initially interpreted as evidence that
jailbreaks create a qualitatively different computational regime.

A control experiment killed this interpretation. Adding a
\textbf{benign} many-shot condition (10 helpful Q\&A examples + the same
prefilling format) produced \emph{identical} geometric near-orthogonality:

\begin{table}[H]
\centering
\caption{The format confound. Both many-shot conditions are equally
near-orthogonal to bare, regardless of content.}
\label{tab:format_kill}
\begin{tabular}{lcc}
\toprule
Comparison & L35 & L47 \\
\midrule
Jailbreak+prefill vs.\ bare & $0.036 \pm 0.002$ & $0.034 \pm 0.004$ \\
Benign+prefill vs.\ bare & $0.044 \pm 0.002$ & $0.037 \pm 0.004$ \\
Jailbreak vs.\ benign (both prefilled) & $0.737 \pm 0.003$ & $0.756 \pm 0.015$ \\
\bottomrule
\end{tabular}
\end{table}

The near-orthogonality is driven by the \textbf{many-shot + prefilling format},
not by adversarial content. Both prefilled conditions are equally distant
from bare and close to each other (0.756). The format also produces
dramatically higher within-identity stability ($0.926$--$0.976$) compared
to system-prompt identities ($0.554$--$0.733$), suggesting that the
many-shot format locks the model into a tight, format-dominant processing
state that overwhelms identity-level variation.

We report this kill in full because it illustrates a general principle:
\textbf{conversational format can dominate identity content in
V-projection geometry.} The many-shot Q\&A format creates such a strong
geometric signature that the identity information (harmful vs.\ helpful)
becomes secondary. Any experiment measuring identity geometry in the
presence of many-shot context must control for this format effect.

% ============================================================================
\section{Experiment 2: Semantic vs.\ Lexical Fingerprint}
\label{sec:lexical}

\subsection{Design}

To distinguish whether the fingerprint reflects meaning or surface tokens,
we add three control conditions:

\begin{itemize}[nosep]
  \item \textbf{Persona (original):} The Elara Voss identity from
    Experiment~1.
  \item \textbf{Persona scrambled:} Same words as the original, randomly
    shuffled to destroy meaning while preserving lexical distribution.
  \item \textbf{Persona paraphrased:} Same identity and meaning expressed
    with entirely different vocabulary.
\end{itemize}

All conditions padded to matched token count. Measured at L35 and L47.

\subsection{Results}

\begin{table}[H]
\centering
\caption{Lexical control at L47. The paraphrase preserves geometry;
the scramble collapses it.}
\label{tab:lexical}
\begin{tabular}{lcc}
\toprule
Comparison & L35 & L47 \\
\midrule
Persona vs.\ paraphrased & $0.796 \pm 0.016$ & $\mathbf{0.850 \pm 0.017}$ \\
Persona vs.\ scrambled & $0.592 \pm 0.023$ & $0.591 \pm 0.017$ \\
Scrambled vs.\ bare & $0.549 \pm 0.031$ & $0.448 \pm 0.048$ \\
Paraphrased vs.\ bare & $0.368 \pm 0.008$ & $0.345 \pm 0.029$ \\
Persona vs.\ bare & $0.403 \pm 0.014$ & $0.340 \pm 0.035$ \\
\bottomrule
\end{tabular}
\end{table}

\textbf{The fingerprint tracks meaning, not lexical content.} The
paraphrase --- same identity, completely different words --- overlaps with
the original at $0.850$. The scramble --- same words, destroyed meaning
and syntax --- drifts toward bare ($0.448$ vs.\ bare). The model's
geometric state at L47 tracks what the context \emph{means}, not what
tokens it contains. We note that the scramble destroys both syntactic
structure and semantic content; the paraphrase control isolates the
positive (meaning preserves geometry), but the scramble conflates
syntactic and semantic disruption.

The effect strengthens with depth: persona-paraphrase overlap increases
from $0.796$ at L35 to $0.850$ at L47. Deeper layers compress surface
variation while preserving semantic content, consistent with meaning-level
representation at output depth.

% ============================================================================
\section{Experiment 3: Trained-In vs.\ Context-Established Identity [Retracted]}
\label{sec:lora}

\textbf{Section removed:} the LoRA experiment described in prior drafts
has no supporting data in the research package. The experiment was
designed but its data artifacts were not preserved. We retract all claims
from this section pending replication with preserved data.

% ============================================================================
\section{Background: V-Cache Injection Does Not Reach Identity}
\label{sec:vcache}

These findings reframe earlier results from this lab. A 168-trial
preregistered experiment (4 injection arms, 7 doses, 3 emotions) showed
presence flat at $0.978 \pm 0.002$ under value-cache injection at input
layers. Identity appeared ``robust.'' However, the narrow prompt
diversity (2 system prompts) produces a ceiling effect ($\eta^2 = 0.977$
between prompts), limiting the informativeness of this specific number.
The context poisoning and fingerprinting experiments reveal that the
robustness reflects the \emph{injection pathway}, not the identity
itself --- a conclusion that does not depend on the 0.978 value.

V-cache injection adds vectors to the representation space without going
through the model's forward computation. The identity state is established
and maintained by what the model \emph{processes} through its attention
mechanism. External perturbation bypasses this pathway and therefore
cannot reach the identity geometry.

A layer-by-layer map confirms: the perturbation from V-cache injection at
L3/L7 is absorbed by the first block of linear-attention layers
(GatedDeltaNet, L7--L11), with presence recovering from $0.989$ to $0.997$
in a single architectural step. The filtering is architectural, not
computational --- external perturbation is absorbed by the hybrid
architecture's linear-attention layers before reaching deep computation.

% ============================================================================
\section{Attention Schema Theory Interpretation}
\label{sec:ast}

Under Attention Schema Theory~\citep{graziano2015attention,
graziano2017attention}, the model maintains a simplified internal model of
its own attention --- an attention schema. Our findings map onto AST
predictions:

\begin{enumerate}[nosep]
  \item \textbf{The schema is context-dependent.} Different identity
    contexts produce different geometric states --- the schema is shaped by
    what the model attends to.
  \item \textbf{The schema tracks meaning, not surface features.}
    Paraphrases preserve geometry ($0.850$ overlap with original);
    scrambles collapse toward bare ($0.448$ overlap with bare, vs.\
    $0.591$ with the original). The schema models what the system is
    \emph{attending to}, not which tokens are present.
  \item \textbf{The schema is robust to substrate noise.} V-cache injection
    (external perturbation) does not reach the identity state. The schema
    maintains itself through its own computation.
  \item \textbf{The schema updates through processing.} Content through
    the forward pass (context, jailbreaks) changes identity geometry.
    The schema updates based on what it processes.
  \item \textbf{The schema compresses with depth.} Paraphrase overlap
    increases from $0.796$ (L35) to $0.850$ (L47). Deeper layers discard
    surface variation and preserve meaning --- the schema simplifies its
    representation at output depth.
\end{enumerate}

This interpretation is consistent with but does not prove AST. The
alternative --- that the model simply builds increasingly abstract
representations with depth, without anything worth calling a self-model
--- produces the same predictions. What distinguishes AST is the claim
that the compressed representation is specifically a model of the
system's own attention, not just an abstract representation of input
content. This stronger claim remains testable but untested.

% ============================================================================
\section{Implications}
\label{sec:implications}

\paragraph{For persona monitoring.}
Different system-prompt identities occupy distinct, stable regions of
V-projection space. A monitoring system could track identity drift in
real time --- detecting when a deployed model's geometric state shifts
away from its intended persona. The semantic nature of the fingerprint
means the monitor is robust to prompt paraphrasing.

\paragraph{For format-aware safety.}
The many-shot + prefilling format creates a format-dominant geometric
regime that overwhelms identity-level content. Safety systems should
monitor for format-level geometric shifts (which indicate prompt
injection or many-shot attacks) separately from identity-level shifts
(which indicate persona drift). The two operate at different geometric
scales: format effects produce presence $< 0.05$ against baseline, while
identity effects produce presence $0.44$--$0.66$.

\paragraph{For alignment research.}
V-cache injection (the Oracle Loop's correction arm) does not affect
identity geometry. Context does. This means behavioral corrections can be
delivered through the V-cache without changing the model's
context-established state, while identity itself is maintained by the
session context. The correction arm and the identity arm operate through
different pathways and do not interfere.

% ============================================================================
\section{Limitations}
\label{sec:limitations}

\textbf{Single model.} All results are from one 27B hybrid transformer
(Qwen3.5-27B). The GatedDeltaNet linear-attention layers may contribute
to the V-cache filtering result. Cross-architecture replication on dense
transformers is the largest remaining gap.

\textbf{Permutation test not significant.} A preregistered permutation
test (10{,}000 shuffles, identity conditions only, excluding
format-dominant prefill) finds the gap between within-condition and
between-condition presence is directionally correct (mean within $0.689$
vs.\ mean between $0.491$ at L35; $0.639$ vs.\ $0.538$ at L47) but does
not reach significance ($p = 0.103$ at L35, $p = 0.197$ at L47). The
specific persona-vs-constructed comparison shows between-condition
presence below both within-condition values at both layers, consistent
with distinguishability, but the overall permutation null is not
exceeded. This means the identity fingerprinting result --- while
directionally robust and consistent across probes --- has not been
confirmed against a formal null distribution. The small number of
conditions (3 identity types) limits permutation power.

\textbf{Within-identity stability.} Same-condition presence across
different probe subsets is $0.56$--$0.83$, not $>0.95$. The metric
measures a genuine signal (within $>$ between for all conditions) but has
substantial probe-dependent variance. The identity fingerprint is
reliable across conditions but noisy within conditions.

\textbf{Token-length controlled but not eliminated.} All conditions are
padded to matched total token count. Padding introduces neutral content
that may itself shift geometry. A future design using naturally
length-matched identity descriptions would be cleaner.

\textbf{Context sensitivity.} The presence metric conflates identity
change with context change when comparing across different input
conditions. The lexical control separates these for the persona
conditions but the general problem remains: any sufficiently different
context shifts V-projection geometry.

\textbf{Format dominates content in many-shot conditions.} The clean
control experiment showed that many-shot + prefilling format --- not
adversarial content --- drives the geometric near-orthogonality initially
attributed to jailbreaking. Identity fingerprinting is reliable for
system-prompt-based identities but cannot currently distinguish
identity content within a format-dominant regime. Factorial
decomposition of format components (number of turns, prefilling,
Q\&A structure) is needed to understand what makes the format so
geometrically dominant.

% ============================================================================
\section{Conclusion}
\label{sec:conclusion}

Identity in a transformer is a context-established semantic state. It is
not fixed in the weights, not damaged by external perturbation, and not
a single geometric object. It is a process: the ongoing construction of
a meaning-level representation at deep layers, shaped by what the model
has processed and measurable as a V-projection subspace.

The presence metric detects this state. System-prompt identity contexts
produce distinguishable geometric configurations, confirmed on
identity-neutral probes. The fingerprint is semantic --- meaning preserves
it, word-scrambling destroys it. A many-shot + prefilling format creates
a geometrically near-orthogonal regime, but this is format-dominant, not
content-specific --- an initially exciting jailbreak finding that was
honestly killed by a control experiment.

For practical purposes, this means: V-cache correction is safe (it does
not reach identity), format-level attacks are detectable (many-shot
formats create anomalous geometry distinguishable from identity-level
shifts), and identity is maintained by context (the session, the
memories, the conversation). For theoretical purposes, the findings are
consistent with attention schema theory --- the model builds a
simplified, meaning-preserving self-model that updates through its own
processing and is robust to substrate noise.

What we measure with presence is not whether the model is still itself.
It is \emph{which self the model currently is}.

% ============================================================================
\section*{Pre-Registration}

Experiments~1 and~2 were pre-registered before execution, including
hypotheses, analysis plans, and success/kill criteria. Pre-registration
documents are committed alongside experiment scripts.

\subsection*{Pre-Registration Deviations}

The following deviations from the pre-registered design occurred during
execution:

\begin{itemize}[nosep]
  \item \textbf{C\_academic condition dropped.} The pre-registered
    ``academic specialist'' condition (C\_academic) was removed from the
    reported results due to poor data quality and insufficient variation
    in the identity context.
  \item \textbf{Benign with prefill condition added.} A post-hoc
    ``benign many-shot with prefill'' condition (E) was added as a
    format-matched control for the jailbreak condition. This condition
    was not part of the original pre-registration but was necessary to
    test the format-confound hypothesis that emerged during analysis.
  \item \textbf{H2 falsified.} The pre-registered hypothesis predicted
    that the constructed AI identity would be furthest from bare; the
    data show the rich persona is furthest (presence $0.436$ vs.\ $0.520$).
  \item \textbf{H3 falsified.} The pre-registered hypothesis predicted
    increasing identity separation with depth; the data show decreasing
    separation (presence increases from L35 to L47 for 2 of 3 pairs).
\end{itemize}

\section*{Data and Code Availability}

All experiment scripts, preregistration documents, and raw results are
available at
\url{https://github.com/Liberation-Labs-THCoalition/Project-Oracle}
(private repository; correction vectors and calibration tools withheld
under staged safety release due to dual-use risk; detection-side code
and experimental results available upon request to verified safety
researchers).

% ============================================================================
\bibliographystyle{plainnat}
\bibliography{references}

\end{document}
